\documentclass[a4paper,12pt]{ctexart} % 使用 ctexart 文档类，原生支持中文

% ==================== 宏包引入 ====================
\usepackage{geometry}       % 页面边距设置
\usepackage{graphicx}       % 插入图片
\usepackage{float}          % 图片位置强制固定
\usepackage{booktabs}       % 三线表
\usepackage{amsmath}        % 数学公式
\usepackage{listings}       % 代码块
\usepackage{xcolor}         % 颜色支持
\usepackage{hyperref}       % 目录跳转和超链接
\usepackage{caption}        % 图表标题格式
\usepackage{subcaption}
\captionsetup[table]{position=bottom,justification=centering,singlelinecheck=false}


% ==================== 页面设置 ====================
\geometry{left=2.5cm, right=2.5cm, top=2.5cm, bottom=2.5cm}

% ==================== 代码高亮设置 ====================
\definecolor{codegreen}{rgb}{0,0.6,0}
\definecolor{codegray}{rgb}{0.5,0.5,0.5}
\definecolor{codepurple}{rgb}{0.58,0,0.82}
\definecolor{backcolour}{rgb}{0.95,0.95,0.92}

\lstdefinestyle{mystyle}{
    backgroundcolor=\color{backcolour},   
    commentstyle=\color{codegreen},
    keywordstyle=\color{magenta},
    numberstyle=\tiny\color{codegray},
    stringstyle=\color{codepurple},
    basicstyle=\ttfamily\footnotesize,
    breakatwhitespace=false,         
    breaklines=true,
    postbreak=\mbox{\textcolor{red}{$\hookrightarrow$}\space},
    captionpos=b,                    
    keepspaces=true,                 
    numbers=left,                    
    numbersep=5pt,                  
    showspaces=false,                
    showstringspaces=false,
    showtabs=false,                  
    tabsize=2,
    columns=flexible
}
\lstset{style=mystyle, language=Python}

% ==================== 文档信息 ====================
\title{\textbf{大模型课程第二次作业}}
\author{姓名：张宗桪 \\ 学号：2023K8009991013 \\ 专业：人工智能}
\date{\today}

% ==================== 正文开始 ====================
\begin{document}

\pagestyle{plain} % 页面顶部不显示标注，仅在底部显示页码

% 封面
\maketitle
\thispagestyle{empty} % 封面不显示页码
\newpage

% 目录
\tableofcontents
\newpage
\section{实验设计}
\subsection{数据集与预处理}
实验采用了混合语料库，包括新华日报官方标注集和THUCNews混合。
\begin{itemize}
    \item \textbf{语料规模}：为了平衡训练效率与效果，最终混合生成的最终语料大小控制在 300MB 左右。
    \item \textbf{预处理逻辑}：去除冗余空白符。为了增强数据随机性。
    \item \textbf{分词器}：使用 BPE算法构建词表，序列截断长度（SEQ\_LEN）设定为 \textbf{512}。
\end{itemize}

\subsection{模型参数配置}
本实验构建了一个 Decoder-only 的模型，具体参数如下表所示：

\begin{table}[H]
    \centering
    \caption{MiniLlama 模型核心参数}
    \begin{tabular}{lc}
        \toprule
        参数名称 & 设定值 \\
        \midrule
        隐藏层维度 (Hidden Dim) & 512 \\
        注意力头数 (N\_Heads) & 8 \\
        网络层数 (N\_Layers) & 6 \\
        FeedForward 维度 & 2048 \\
        激活函数 & GELU \\
        \bottomrule
    \end{tabular}
\end{table}

\subsection{训练策略}
\begin{itemize}
    \item \textbf{优化器}：采用 AdamW 优化器，学习率设为 1e-4，并配合 \texttt{weight\_decay=0.01}。
    \item \textbf{调度器}：使用余弦退火策略动态调整学习率。
    \item \textbf{精度与效率}：开启了 \texttt{bfloat16} 混合精度训练，并设定 Batch Size 为 32。为了防止训练不稳定，加入了阈值为 1.0 的梯度裁剪。
\end{itemize}

\section{实验结果}

\subsection{Loss 曲线分析}

\begin{figure}[H]
    \centering
    \includegraphics[width=0.95\textwidth]{loss_curve_20260418_152103.png}
    \caption{训练过程 Loss 下降曲线}
\end{figure}


\begin{enumerate}
    \item \textbf{收敛情况}：模型 Loss 从初始的 10.2 左右快速下降，最终在约 15,000 个 Step 后稳定在 5.0 附近。Epoch 平均 Loss 曲线表现极其平稳，证明了训练过程中的超参数设置非常有效，确保了训练的稳定性。
    \item \textbf{语言规律学习}：在训练的前 10\% 阶段，Loss 下降最为显著，这对应了模型从完全随机状态过渡到掌握基础词汇概率分布的过程。
    \item \textbf{拟合程度}：随着训练进入第三个 Epoch，Loss 曲线逐渐变平
\end{enumerate}

\subsection{文本续写展示}
使用训练得到的权重进行推理测试，给定前缀后的生成结果如下：

\begin{flushleft}
\textbf{Prompt：} 我是一个粉刷匠。\\[0.5em]
\textbf{Continuation：}
\begin{quote}
我 是一个 粉 刷 匠 。 虽然 只有 没有 小 贩 卖 , 所以 将 她 就 能 将 去 找 出 , 而且 她 对 她 带 走 我 很 小 了 自己的 小 米 , 但 这 不 知道 了 , 所以 她 要 这样 的一 局 把 自己 , 他 因为 她 就 要 对 爱 情 , 所以 她 决定 找 来 找 著 她 , 但是 她 一直 找 出 色的 孩 子的 生活 , 她 更 像 与 她 相 逢 , 但 她 表示 她的 生活 , 并且 自己 是 和 亲 情 , 是 这 么 对 这 么 。 但 她 不 起 , 她在 她 想 就 能 找到 自己的 男人 , 因为 她 也 就 喜欢 和 她 的 她 生活 。 然而
\end{quote}
\end{flushleft}


\begin{flushleft}
\textbf{Prompt：}中国科学院大学是\\[0.5em]
\textbf{Continuation：}
\begin{quote}
中国科学院 大学 是 该校 的 博士 。 大学 的 学部 学系 中 课程 也是 艺术 运动 会的 学生 , 一个 大学 学生 在 学 上 获得 教师 。   大学 校园 的 课程 和 学 和 课程 的 教学 是 , 学院 学生 学 、 教师 和 学 等 教育 课程 , 与 高中 生 之间 , 在 学院 时期 , 学院 课程 管理 , 学生 与 讲 学 。   校 史  *   中 立 大学 短期大学 与 研究生 学院 , 学校 学习 。 该校 为 英语 大学 , 校 生 、 体育 学 、 学 硕士 , 教育 学 博士 生 , 研究生 学位 , 校 生 导师 。   学生   *   中 立 大学 , 数学 学 和 讲 学 , 课程 。   *
\end{quote}
\end{flushleft}

\section{结果分析与见解}

\subsection{模型收敛性分析}
从 Loss 曲线的表现来看，模型基本达到了收敛状态。在训练过程中没有出现明显的梯度爆炸或震荡，这主要归功于模型初始化时的参数设置以及 AdamW 优化器的稳定性。虽然最终 Loss 并没有降到极低，但考虑到 300MB 的语料量和仅 18M 左右的参数量，当前的收敛效果符合预期。

\subsection{生成能力评估}
针对上述生成结果，结合训练过程中的 Loss 曲线，得出以下结论：
\begin{enumerate}
    \item \textbf{语言习得能力}：模型已经学会了基本的中文标点使用和词汇搭配。即使在语义不通的情况下，生成的句子结构依然保持了较高的完整性，证明了自回归预训练对捕获统计规律的有效性。
    \item \textbf{逻辑与事实分析}：由于模型参数量仅为 18M 且训练轮数有限，生成内容表现出显著的“逻辑碎片化”特征。模型无法维持长程语义的自洽性，且在事实性陈述（如学校定义）上表现出明显的幻觉现象。

\end{enumerate}

\end{document}